% =====================================================================
% Section VIII -- A Catalog of Measurement Artifacts
% SpecAsync-UVM manuscript, draft 1 (2026-11-06)
% Source: ARTIFACT_CATALOG.md (31 entries), RESULTS_FACTS.md Section VIII
% (family counts and how-caught column), T-CAT (paper/tables/t_cat.tex).
% Entry numbers are the catalog's. The "103 days" for #13 is from the
% catalog text; confirmed in C15 (2026-06-13 to 2026-09-24).
% =====================================================================

\section{A Catalog of Measurement Artifacts}
\label{sec:pitfalls}

Over the course of this study we recorded 31 artifacts: errors in measurement,
analysis or procedure, each with the false conclusion it produced or would have
produced. We report them for two reasons.
Several of them shaped the results of Sections~\ref{sec:spec}
and~\ref{sec:threshold}, including the reversal of this study's earlier
conclusion, so a reader needs them to judge those results. And most are not
specific to our mechanism: they can arise in any measurement of a GPU driver.

\subsection{Families}

We group the entries by the mechanism of the error, not by its consequence
(Table~\ref{tab:families}).

\begin{table}[t]
\centering
\footnotesize
\caption{Catalog entries by family.}
\label{tab:families}
\begin{tabular}{@{}l l r@{}}
\toprule
Family & Entries & $n$\\
\midrule
Aggregation-level mismatch & 7, 8, 9, 13, 16, 17, 18, 25, 27 & 9\\
Validation that could not fail & 1, 2, 12, 14, 15, 31 & 6\\
Instrument cost and blindness & 4, 5, 19, 20, 21 & 5\\
Operational hazards & 3, 10, 11, 22, 24, 28, 29, 30 & 8\\
Record integrity & 23, 26 & 2\\
Uncategorised (a genuine outlier) & 6 & 1\\
\bottomrule
\end{tabular}
\end{table}

\paragraph{Aggregation-level mismatch}
A quantity is recorded at one level and consumed at another: per batch and
per fault, per trial and per run, a share of one window applied to another. The
most consequential entry in the catalog is of this kind (\#13). Our first
oracle recorded faults once per batch and replayed them once per fault, so it
ran ahead of the application and wrapped. Every oracle result for 103 days was
read as the result of a maximally favourable predictor, and the conclusion
that ``even oracle prediction yields no benefit'' was the paper's thesis. The
others in the family are smaller versions of the same error: capped ring
counts read as totals (\#7), rates computed over the wrong denominator (\#8,
\#16), a ceiling expressed in the units of a different window (\#9, \#25), and
a multiplicity correction hard-coded for a family of 16 and reused for others
(\#27, Section~\ref{sec:protocol}).

\paragraph{Validation that could not fail}
A check passes, but it would have passed whatever the truth was. The first
fix to the oracle was validated by an improvement, without an absolute
ceiling against which to judge it (\#15). A positive control that could not have
failed was read as showing that the fix, its hit accounting and the width path
were sound (\#14). A large speedup under oversubscription,
measured against a crippled baseline, appeared to reverse the paper's
conclusion (\#12). The newest entry is
of this kind too (\#31). We validated GPU memory in use as a measure of
migrated volume on HMM memory, where the driver allocates 4\,KB pages, and then
used it on managed memory, where the driver allocates a whole 2\,MB chunk per
block on the first fault. There it reads allocation, not migration, and an
expectation scored on it held (Section~\ref{sec:thr-hmm}).

\paragraph{Instrument cost and blindness}
An instrument changes what it measures, or cannot see it. The clearest case is
our oracle (\#21): its lookup ran on the fault-servicing thread under a global
lock, so an instrument meant to bound speculation's benefit was paying most of
speculation's per-fault cost (Section~\ref{sec:spec-cost}). Another counter,
named \texttt{spec\_hits}, counted faults on pages predicted within a time
window and was read as prediction success; a second counter, measuring whether
the worker actually finished first, disagreed with it (\#20). Two entries
concern our own ring buffers, read at a drifting byte offset (\#5) and in slot
order rather than time order (\#19).

\paragraph{Operational hazards}
These are errors of procedure: results taken with address randomisation
disabled read as general platform behaviour (\#10), a host memory leak (\#11), a kernel-log diff
by line count that could miss lines (\#24), a driver patch silently replaced by a Git LFS
pointer (\#28), a session ended by our own clean-up command (\#29), and an
oversubscribed run sized against GPU memory but not host memory (\#30). The one
that changed our development protocol is \#22: an instrumentation counter called
a driver helper that returns a null pointer when the GPU has no state in a
block, guarded only by a non-fatal assertion, and the kernel oopsed
(Section~\ref{sec:design-safety}).

\paragraph{Record integrity}
Two entries concern our records rather than our measurements: a platform table
that silently fell behind a driver upgrade (\#23), and a status log treated as
a primary record of timing (\#26). Section~\ref{sec:protocol} explains why we
establish the order of events by version-control ancestry, not by the clocks in
those records.

\subsection{How the Errors Were Caught}

Each entry records how it was found. Of the nine aggregation-level errors, eight
were caught only by a later, independent pass: a re-investigation, a source
audit, or the author of a later gate. Seven of them had already reached a
reported figure or conclusion. One (\#25) was caught inside the gate that made
it, by a pre-registered prediction test: E0.9b's addendum predicted, from pilot
runs, the interval in which its result should fall, and the result fell outside
it. Across all families, four errors were caught inside their own gate, and all
four by a check written before the measurement: a pre-registered test (\#10,
\#21, \#25) or a pre-registered smoke pass (\#30).

We state this pattern descriptively. The samples are small, two or three entries
per catching mechanism, and the counts are biased: an error caught inside its gate may
never have been recorded as an artifact at all, so in-gate catches are
under-counted by construction. Two lessons nevertheless shaped the protocol of
Section~\ref{sec:protocol}. A check that cannot fail is not a check, hence the
positive controls that must also reject a perturbed copy. And an error of
aggregation level is invisible from inside the analysis that makes it, because
every step is arithmetically correct; what catches it is a second, independent
path to the same number.

The full catalog, with the false conclusion and the means of detection for each
entry, is in the supplementary material (Section S3).
